\pageheading{PRIME SIEVES\; /\; MULTIPLE PARAMETERS}{Overlap is part of the arithmetic}
For a finite set $F$ of distinct primes, define
\[
\nu_{F,\boldsymbol\lambda}
=\sum_{n\ge1}\mu(n)^2n^{-4}
\prod_{p\in F}(1+\lambda_p\mathbf1_{p\mid n})\delta_{x_n},
\quad x_n=(1-n^{-2})^2.
\]
For real strengths, the exact positivity domain is
\[
\boxed{\lambda_p\ge-1\quad\text{for every }p\in F.}
\]
Sufficiency is termwise positivity. If one strength is below $-1$, the
atom $n=p$ isolates its negative factor. The other strengths do not occur
at that atom. At every simultaneous deletion endpoint, primes outside
$F$ still supply infinitely many positive nodes, so all Hankel matrices
remain strictly positive definite.

The exact coordinate formula is
\[
M_j(\boldsymbol\lambda)=\sum_{r=0}^{2j}(-1)^r\binom{2j}{r}
\frac{\zeta(4+2r)}{\zeta(8+4r)}
\prod_{p\in F}\left(1+\frac{\lambda_p}{p^{4+2r}+1}\right).
\]
Coefficients of the monic families lie in
$\mathbb Q(\pi^{-2},\lambda_p:p\in F)$.
For two primes, the overlap term is visible:
\[
M_j=M_j(0,0)+\lambda_pU_{p,j}+\lambda_qU_{q,j}
+\lambda_p\lambda_qU_{pq,j}.
\]
Here $U_{pq,j}>0$ is the $pq$-divisible submeasure moment. The family
does not reduce to adding separate prime responses without their overlap.

\pageheading{PRIME SIEVES\; /\; FINITE TESTS}{Moment positivity is a weaker requirement}
\fig{12_multi_prime_domain}{The exact positive-measure domain and the
order-one moment test have different parameter geometry. Regions passing
a finite test outside the safe orthant still contain negative atoms.}
For primes two and three,
\[
M_0=105\pi^{-4}(1+\lambda_2/17)(1+\lambda_3/82).
\]
Both $(0,0)$ and $(-34,-164)$ have positive total mass. Their midpoint
$(-17,-82)$ has zero mass. Thus the parameter set passing the order-one
positive-definiteness test is nonconvex. This particular midpoint argument
does not establish nonconvexity of the semidefinite test, which permits zero.
For that test, take the two points $(0,0)$ and $(-34,-246)$.
The normalized masses $M_0/(105\pi^{-4})$ there are $1$ and $2$.
Their convex combination with weights $3/5$ and $2/5$ is
$(-68/5,-492/5)$, where the normalized mass is $-1/25$.
Thus the order-one semidefinite feasible set is also nonconvex,
by a separate counterexample.

\claim{A separate additive family has different geometry}
For weights $1+\sum_{p\in F}\gamma_p\mathbf1_{p\mid n}$, positivity is
equivalent to
\[
1+\sum_{\gamma_p<0}\gamma_p\ge0.
\]
Every prime subset appears at its squarefree product, proving necessity.
Checking each atom proves sufficiency. Its finite Hankel matrices are affine
in $\gamma$, so the semidefinite feasible sets are convex spectrahedra.
The multiplicative family has genuine overlap terms and polynomial parameter
dependence. These are two distinct deformations, with distinct domains.

\pageheading{PRIME SIEVES\; /\; CONSTRUCTIVE WITNESS}{Every negative prime atom has a finite detector}
For an integer $a\ge2$, let
\[
L_a(x)=\prod_{n=1}^{a-1}\frac{x-x_n}{x_a-x_n},\quad
C_a=L_a(1),\quad r_a=\frac{1-x_{a+1}}{1-x_a}<1,
\]
\[
f_{a,m}(x)=L_a(x)\left(\frac{1-x}{1-x_a}\right)^m.
\]
This polynomial kills every lower node, equals one at $x_a$, and has
absolute value at most $C_ar_a^m$ at every higher node.
\fig{13_signed_detector}{A rational detector for $p=3$, $\lambda_3=-1.1$,
$\lambda_2=1$. The first exact bound proving negativity is $m=5$, so a
Hankel matrix of order eight detects the signed defect.}
Set $C_F=\prod_{q\in F}\max(1,|1+\lambda_q|)$. The summable tail gives
\[
\int f_{p,m}^2\dd\nu
\le\frac{1+\lambda_p}{p^4}+
\frac{C_FC_p^2}{3p^3}r_p^{2m}.
\]
Hence the quadratic form is negative whenever
$pC_FC_p^2r_p^{2m}/3<-(1+\lambda_p)$.
For rational strengths this is an exact rational certificate. It needs no
numerical eigenvalue calculation. A depth tilt $x^k$ replaces the negative
atom contribution by $(1+\lambda_p)x_p^k/p^4$.

\pageheading{PRIME SIEVES\; /\; FIBERS}{Real singularities survive multiple deletions}
Fix all strengths except $\lambda_p$, with the others at least $-1$.
The remaining moment matrix has the form $A_m+(1+\lambda_p)B_m$, where
$A_m$ is the $p$-free part and $B_m$ the $p$-divisible part. Both are
positive definite, even at other deletion endpoints: primes outside $F$
supply the first support, and integers $p\ell$ the second.

Thus every one-variable determinant fiber has exact degree $m$ and all
roots real below $-1$. Its nearest margin satisfies
\[
\delta_m=\min_{\substack{f\in\mathbb R[x]\setminus\{0\}\\\deg f<m}}
\frac{\int f^2\dd\nu_{p\text{-free}}}
{\int f^2\dd\nu_{p\text{-divisible}}},\qquad\delta_m\downarrow0.
\]
The detector yields, for $C_\kappa=\prod_{q\ne p}\max(1,1+\lambda_q)$,
\[
\delta_{p+m}\le\frac{pC_\kappa C_p^2}{3}r_p^{2m}.
\]
The denominator retains the atom $p$; all lower nodes in the numerator are
killed. This extends the prime-two bound to every prime and to simultaneous
positive or deletion strengths. It does not assert simple roots or global
multivariate real stability.

If at least two positive-domain strengths are nonzero, the complete monic
sequence cannot equal a one-prime sequence on the same variable and nodes.
Compact moment determinacy equates normalized measures; the fixed mass at
zero fixes their scale; atom weights at the individual primes then read off
every strength. An arbitrary affine change of variable is outside this separator.
